% Part: proof-theory
% Chapter: proof-search
% Section: introduction

\documentclass[../../../include/open-logic-section]{subfiles}

\begin{document}

\olfileid{pt}{ps}{int}

\olsection{પરિચય}

\Log{G3c} અથવા
\Log{N1i} જેવા
!!{proof}-તંત્રોના
સ્વયંસિદ્ધો અને
અનુમાન-નિયમો
નક્કી કરે છે
કે સિક્વન્ટો
અથવા !!{formula}sનાં
કયાં વૃક્ષો
યોગ્ય !!{proof}s
ગણાય. આવું
સિક્વન્ટો અથવા
!!{formula}sનું
વૃક્ષ અપાયું
હોય (દરેક
પગલે કયો નિયમ
લાગ્યો છે
અથવા ધારણાઓ
અને તેમને
મુક્ત કરતા
નિયમોનાં
ચિહ્નો જેવી
વધારાની માહિતી
સાથે પણ),
તો સ્વયંસિદ્ધો
અને નિયમોની
વ્યાખ્યાઓથી
તે યોગ્ય
!!{proof} છે
કે નહીં તે
તપાસી શકીએ.
પણ અપાયેલા
સિક્વન્ટની
અથવા અપાયેલી
ધારણાઓમાંથી
અપાયેલા
!!{formula}ની
!!a{proof}
\emph{શોધવી}
અલગ અને
સામાન્ય રીતે
વધુ મુશ્કેલ
સમસ્યા છે.
આ \emph{સાબિતી-શોધ}
ની સમસ્યા છે.

સાબિતી-શોધ માટે
ખૂબ સાદી
રીત છે.
!!{formula}sને
સીમિત વર્ણમાળાના
ચિહ્નોના ક્રમ
તરીકે વિચારી
શકાય. સિક્વન્ટો
અને સિક્વન્ટો
અથવા
!!{formula}sનાં
વૃક્ષોને પણ
!!{formula}sના
ક્રમ તરીકે
વિચારી શકાય
(વૃક્ષની શાખાઓ
બતાવવા ખાસ
કૌંસ વાપરીને).
સ્વયંસિદ્ધો
અને નિયમો
વડે ચિહ્નોનો
ક્રમ યોગ્ય
!!{proof}
દર્શાવે છે
કે નહીં તે
યાંત્રિક રીતે
તપાસી શકાય.
તેથી \emph{બધી
શક્ય} યોગ્ય
!!{proof}s
યાંત્રિક રીતે
બનાવી શકાય:
પ્રથમ !!{proof}s
દર્શાવતી
વર્ણમાળાના
ચિહ્નોના બધા
ક્રમ બનાવીએ
અને દરેક
ઉમેદવાર યોગ્ય
!!{proof} છે
કે નહીં તપાસીએ.
આ સાદી
સાબિતી-શોધ
એવી છે કે
અપાયેલા સિક્વન્ટની
!!{proof} મળે
ત્યાં સુધી
બધી યોગ્ય
!!{proof}s
બનાવતાં રહીએ
(અથવા અપાયેલી
ધારણાઓમાંથી
ખુલ્લી રહે
તેવી ધારણાઓ
સાથે અપાયેલા
!!{formula}ની
સાબિતી મળે
ત્યાં સુધી).

વધુ સારી
રીત એ છે
કે હાથથી
!!a{proof}
શોધવા વપરાતી
સહજ રીત
અપનાવીએ:
અંતિમ સિક્વન્ટથી
શરૂ કરી
!!a{formula}
પસંદ કરીએ
અને અનુમાન-નિયમ
``ઊલટી દિશામાં''
લગાવી એક
અથવા વધુ
આધારવિધાન
બનાવીએ; તેમની
!!{proof}s
હોય, તો
છેલ્લા અનુમાન
સાથે જોડીને
અપાયેલા સિક્વન્ટની
!!a{proof}
મળે. સ્વાભાવિક
નિગમન-તંત્રમાં
!!a{proof}
શોધવા આવી
જ, જોકે
વધુ જટિલ,
રીત વપરાય.

પણ આ રીત
અચૂક નથી:
ખોટો નિયમ
અથવા
!!{formula}sનો
ખોટો ક્રમ
પસંદ કરો,
તો આગળ
માર્ગ ન
રહે. તે
સંપૂર્ણ નિર્દિષ્ટ
પણ નથી:
દરેક તબક્કે
પસંદ કરવા
એકથી વધુ
!!{formula}
અથવા ઊલટી
દિશામાં
લગાવવા એકથી
વધુ નિયમ
હોઈ શકે.
\emph{સાબિતી-શોધ
અલ્ગોરિધમ}
એવી સંપૂર્ણ
નિર્દિષ્ટ
પ્રક્રિયા છે
જે !!a{proof}
હોય તો
તે શોધે,
એટલે અચૂક
હોય.

કેટલાંક
!!{proof}-તંત્રોમાં
સાબિતી-શોધના
અલ્ગોરિધમ
બીજાં કરતાં
સરળ બને છે.
સિક્વન્ટ-તંત્રોમાં
!!a{formula}નો
!!{main operator}
અને તે ડાબે
છે કે જમણે
એ ઊલટી
દિશામાં લગાવવાનો
નિયમ નક્કી
કરે છે.
દાખલા તરીકે,
$!A \land
!B$ જમણી
બાજુ હોય એવા
$\Gamma \Sequent \Delta, !A \land !B$ની
!!a{proof}
શોધવી હોય,
તો
\RightR{\land}
ઊલટી દિશામાં
લગાવી
$\Gamma \Sequent \Delta, !A$
અને
$\Gamma \Sequent \Delta, !B$
એ બે અનુરૂપ
આધારવિધાન
મેળવવાં.
પણ તંત્ર
\Log{G1c}
હોય તો
સંકોચનને લીધે
વિકલ્પો વધે
છે: તમે
\RightR{\Contraction}
ઊલટી દિશામાં
લગાવી તેના
સ્થાને
$\Gamma \Sequent \Delta, !A \land
!B, !A \land !B$
આધારવિધાન
પણ બનાવી
શકો. પછી
$\Gamma \Sequent
\Delta, !A \land !B, !A \land !B, !A \land !B$
અને એમ
આગળ.
\Log{G2c}માં
તો વધુ
મુશ્કેલ છે.
\RightR{\land}
ઊલટી દિશામાં
લગાવવા
$\Gamma_1$,
$\Gamma_2$
એવા અનુમાનવા
પડે કે
$\Gamma
= \Gamma_1 \cup \Gamma_2$;
એ જ રીતે
$\Delta_1$,
$\Delta_2$
એવા કે
$\Delta
= \Delta_1 \cup \Delta_2$.
પછી
$\Gamma_1 \Sequent \Delta_1, !A$
અને
$\Gamma_2 \Sequent \Delta_2, !B$
આધારવિધાનો
અજમાવવાં
પડે. ખોટું
અનુમાન પાછળથી
બંધ માર્ગે
લઈ જાય,
તો શરૂઆતમાં
પાછા ફરી
અલગ
$\Gamma_1$,
$\Gamma_2$,
$\Delta_1$,
$\Delta_2$
પસંદ કરવા
પડે.
!!{formula}
$!A \lor !B$
હોય તો
\RightR{\lor}
માટેના બે
વિકલ્પોમાંથી
એક પસંદ
કરી
$\Gamma \Sequent \Delta,
!A$
અથવા
$\Gamma \Sequent \Delta, !B$
આધારવિધાન
બનાવવું પડે.
ખોટી પસંદગી
ફરી પાછળ
ફરવા મજબૂર
કરે છે.

\Log{G3c}માં
આ સમસ્યાઓ
નથી. તેમાં
બંધારણીય
નિયમો નથી,
અને !!{main operator}
તથા !!{formula}
કઈ બાજુ છે
તે પરથી
અનુમાન-નિયમ
નક્કી થાય
છે. તેથી
\Log{G1c} કે
\Log{G2c} કરતાં
\Log{G3c}
સાબિતી-શોધ
માટે વધુ
અનુકૂળ છે.
સફળ શોધ
\Log{G3c}માં
!!a{proof}
આપે છે,
અને પછી
તે !!{proof}ને
\Log{G1c}
અથવા
\Log{G2c}
(અથવા
\Log{N1c})
માં રૂપાંતરિત
કરી શકાય.
આમ \Log{G3c}
માટે સાબિતી-શોધ
અલ્ગોરિધમ
અને
\Log{G3c}ની
!!{proof}sને
બીજાં તંત્રોની
!!{proof}sમાં
રૂપાંતરિત
કરતા અલ્ગોરિધમ
હોય, તો
તે તંત્રો
માટે પણ
સાબિતી-શોધ
ઉકેલાય છે.

સાબિતી-શોધ
અલ્ગોરિધમ
અચૂક છે
તે કેવી
રીતે જાણીએ?
જો અલ્ગોરિધમ
!!a{proof}
શોધવામાં
નિષ્ફળ જાય,
તો કોઈ
!!{proof}
હોઈ જ
ન શકે
તે બતાવવું
પડે. ખરાબ
અલ્ગોરિધમ
પસંદ કરીએ
તો સાબિતી
હોવા છતાં
તે શોધવામાં
નિષ્ફળ જઈ
શકે. સાદા
અલ્ગોરિધમને
આ મુશ્કેલી
નથી, કારણ
કે તે
બધી શક્ય
!!{proof}s
ચકાસે છે.
પણ સાબિતી-શોધ
અલ્ગોરિધમ
કાર્યક્ષમ
અને ઓછામાં
ઓછો પ્રયત્ન
કરતો હોવો
જોઈએ.

સાબિતી-શોધ
અલ્ગોરિધમ
અચૂક છે
તે બતાવવાનો
મુખ્ય વિચાર
આ છે: અલ્ગોરિધમ
નિષ્ફળ જાય,
તો તેની
નિષ્ફળતાની
રીત પરથી
સિક્વન્ટની
!!a{proof}
શક્ય નથી
તે બતાવીએ.
તે માટે
સિક્વન્ટ
પ્રામાણ્ય નથી
એવું બતાવીએ.
શાસ્ત્રીય
પ્રથમ-ક્રમ
તર્કમાં
$\Gamma \Sequent \Delta$
પ્રામાણ્ય ત્યારે
છે જ્યારે
દરેક
!!{structure}
$\Gamma$ના
ઓછામાં ઓછા
એક !!{formula}ને
અસત્ય અથવા
$\Delta$ના
ઓછામાં ઓછા
એક !!{formula}ને
સત્ય બનાવે.
\Log{G3c}
\emph{યથાર્થ}
છે: તે ફક્ત
પ્રામાણ્ય
સિક્વન્ટો
!!{prove}s
કરે છે.
આ !!{proof}s
પર આગમનથી
સ્થાપિત થાય
છે: સ્વયંસિદ્ધો
સ્પષ્ટ રીતે
પ્રામાણ્ય છે,
અને નિયમનાં
આધારવિધાનો
પ્રામાણ્ય હોય
તો તેનું
નિષ્કર્ષ પણ
પ્રામાણ્ય છે.
સાબિતી-શોધ
અલ્ગોરિધમ
અચૂક છે
તે માટે
બતાવીએ કે
$\Gamma \Sequent \Delta$ની
!!a{proof}
શોધ નિષ્ફળ
જાય, તો
!!a{structure}
એવી વ્યાખ્યાયિત
કરી શકાય
જેમાં
$\Gamma$નાં
બધાં
!!{formula}s
સત્ય અને
$\Delta$નાં
બધાં
!!{formula}s
અસત્ય હોય.
આથી
\Log{G3c}
\emph{પૂર્ણ}
પણ છે:
$\Gamma \Sequent \Delta$
પ્રામાણ્ય
હોય, તો
તેની
\Log{G3c}-!!{proof}
હોય છે.
દરેક નિષ્ફળ
શોધ
$\Gamma \Sequent \Delta$
અપ્રામાણ્ય
છે તે બતાવે
છે, એટલે
પ્રામાણ્ય
સિક્વન્ટની
દરેક શોધ
સફળ થાય.

\end{document}
